\section{七、项目技术的支撑}

\subsection{1.监测技术升级的现实需求}

黑土地是保障国家粮食安全的重要战略资源，其退化过程具有缓变、隐蔽、空间异质性强等特点，只有依靠高频次、高精度、可重复的观测手段，才能及时捕捉退化趋势并支撑保护性耕作决策。长期以来，东北黑土区的监测工作主要依赖传统手段，与常态化、精细化的保护需求之间存在明显差距，具体表现为以下几个方面。

\begin{itemize}
    \item \textbf{人工巡查效率低、覆盖范围有限。}以人工实地采样与巡田为主的方式，受人力、车辆、道路条件与作业季节的制约，单次可覆盖的面积十分有限，难以在作物生长关键期内完成大范围同步普查；同时，人工观测结果依赖经验判断，不同人员在判读口径上难以完全统一，数据的可比性与可追溯性不足。
    \item \textbf{常规卫星遥感分辨率不足。}中低分辨率遥感产品虽然具备大范围覆盖能力，但其像元尺寸往往大于田间地块的实际尺度，混合像元效应显著，难以识别地块尺度的退化细节，如侵蚀沟发育、耕作层变薄、地表裸露程度变化等，往往只能反映区域平均状态，无法支撑逐地块的差异化管理。
    \item \textbf{国外高分辨率数据获取成本高、周期长。}国外亚米级遥感数据在采购费用、重访安排与数据交付环节均存在较多限制，成本与周期难以匹配常态化监测的需要，且数据获取的稳定性与连续性存在不确定性，不利于形成长时序、可对比的监测档案。
\end{itemize}

综合以上短板可以看出，黑土地保护工作迫切需要一套\textbf{国产化、高精度、大范围}的遥感监测方案：以自主可控的数据源为基础，摆脱对外部数据供给的依赖；以亚米级成像能力为支撑，将观测尺度下沉到地块级；以稳定的重访与交付节奏为保障，形成可持续运行的数据链条，从而为黑土地保护提供稳定、可靠的数据支撑。本项目选择国产高分二号卫星作为主数据源，正是基于上述需求做出的技术判断。

\subsection{2.GF-2卫星的核心技术优势}

高分二号（GF-2）卫星是我国自主研发的首颗亚米级民用光学遥感卫星，专为高分辨率对地观测设计，其技术参数与本项目的监测需求高度契合：一方面，其空间分辨率足以支撑地块尺度的地物判读；另一方面，其幅宽与重访能力可以满足区域尺度的周期性覆盖要求。两者结合，使“既看得清、又看得全、还看得勤”成为可能。

从监测需求出发，GF-2的关键技术参数与本项目形成逐条对应的支撑关系：空间分辨率对应识别精度需求，决定地块边界与侵蚀沟能否被分辨；成像幅宽对应覆盖效率需求，决定区域尺度监测所需的成像次数；轨道与重访能力对应时间频次需求，支撑“周级”动态监测持续运行；载荷成熟度与产品供给稳定性对应工程可靠需求，决定长期业务化运行的风险水平。上述参数与数据保障条件彼此配合，共同支撑起从影像获取到变化检测的完整链路。表~\ref{tab:gf2}汇总了 GF-2 的主要技术参数及其对本项目的支撑作用。

\begin{table}[htbp]
    \centering
    \caption{高分二号（GF-2）卫星主要技术参数及其对本项目的支撑}
    \label{tab:gf2}
    \small
    \begin{tabular}{|p{2.6cm}|p{5.6cm}|p{5.8cm}|}
        \hline
        \textbf{参数类别} & \textbf{技术指标} & \textbf{对本项目的支撑作用} \\
        \hline
        发射与轨道 & 2014 年 8 月 19 日发射；太阳同步回归轨道，轨道高度 631 km，回归周期 69 天 & 成像光照条件稳定、重访规律，保障多时相影像可比 \\
        \hline
        空间分辨率 & 星下点全色 0.8 m、多光谱 3.2 m & 亚米级成像可清晰分辨地块边界、植被覆盖度与侵蚀沟等细部特征 \\
        \hline
        成像幅宽 & 45 km（两台相机组合） & 单景覆盖范围大，减少区域监测所需的影像拼接与重访次数 \\
        \hline
        侧摆与重访 & 侧摆能力 $\pm35^\circ$，侧摆条件下重访周期约 5 天 & 支撑“周级”动态监测的频次要求 \\
        \hline
        载荷谱段 & 全色 0.45--0.90 $\mu$m；多光谱：蓝 0.45--0.52、绿 0.52--0.59、红 0.63--0.69、近红外 0.77--0.89 $\mu$m & 覆盖植被与土壤诊断光谱区间，支持植被指数计算与地力指标反演 \\
        \hline
        定位精度 & 无地面控制点时约 50 m（CE90） & 保障多时相影像几何配准精度，减少配准误差引起的伪变化 \\
        \hline
        数据供给 & 国产民用高分辨率卫星，产品成熟稳定 & 数据来源自主可控，保障长时序监测档案的连续性 \\
        \hline
    \end{tabular}
    \par\vspace{2pt}
    {\footnotesize 注：表中参数依据高分二号卫星官方公布数据整理。}
\end{table}

\begin{figure}[htbp]
    \centering
    \fbox{\parbox[c][4.5cm][c]{0.8\textwidth}{\centering 【图片空位：image6.png】高分二号（GF-2）卫星技术参数示意}}
    \caption{高分二号（GF-2）卫星技术参数说明（一）}
    \label{fig:4a}
\end{figure}

\begin{figure}[htbp]
    \centering
    \fbox{\parbox[c][4.5cm][c]{0.8\textwidth}{\centering 【图片空位：image7.jpeg】高分二号（GF-2）卫星成像能力与参数说明}}
    \caption{高分二号（GF-2）卫星技术参数说明（二）}
    \label{fig:4b}
\end{figure}

\subsubsection{（1）空间分辨率：从“看得见”到“看得清”}

全色波段分辨率达0.8米，多光谱波段分辨率为3.2米，属于亚米级高清成像水平。这一分辨率水平意味着单个像元对应的地面尺寸小于一米，已明显优于田间常见地物的尺度阈值：地块边界、田埂与道路走向在影像上可清晰分辨，植被覆盖度的疏密差异可通过纹理与光谱组合加以区分，侵蚀沟等线状微地貌也能被有效识别。对于黑土地精细化监测而言，空间分辨率直接决定了退化指标能否被观测到——只有将观测尺度下沉至地块一级，才能对黑土区地块边界、植被覆盖度、侵蚀沟等细微地物特征进行稳定判读，进而满足黑土地精细化监测的精度要求。全色与多光谱的波段组合，则在保证几何细节的同时保留了地物光谱信息，为后续变化检测模型提供了信息量更充分的数据基础。

亚米级分辨率的价值，在地块边界与侵蚀沟两类目标上体现得尤为直接。地块边界多由田埂、道路、沟渠等窄线状地物构成，宽度往往不足数米，在中低分辨率影像上极易被混合像元淹没，而在高分辨率影像上则呈现为连续清晰的线状特征，使地块几何轮廓得以准确勾画，这是开展逐地块统计与差异化管理的前提。侵蚀沟属于典型的缓变型线状微地貌，发育初期宽度较小、与周边耕地光谱差异有限，只有在足够精细的空间尺度上才能被稳定识别；将其纳入常态化监测对象，有助于在沟蚀尚处初期时发现问题，避免其扩展造成耕作层不可逆损失。

\subsubsection{（2）成像幅宽：从“看得清”到“看得全”}

单次成像幅宽达45公里，可实现大面积区域的同步覆盖。幅宽是决定监测效率的关键参数：幅宽越大，单轨可覆盖的地面范围越广，拼接与配准的工作量越小，同一时相内影像的辐射与几何一致性也越好。相较于幅宽偏窄的高分辨率卫星，GF-2能够在更少的成像次数内完成同等面积的观测，成像效率远超国外同类卫星。这一特性与松嫩平原等黑土集中连片分布区的空间形态高度匹配，能够高效完成东北黑土区大范围、周期性的动态监测任务，为“周级”动态监测的频次目标提供了工程上可行的时间余量。

\subsubsection{（3）轨道与成像能力：从“看得全”到“看得懂”}

卫星采用太阳同步轨道，轨道高度约631公里，可稳定获取高重访率的观测数据。太阳同步轨道使卫星在相近的地方时过境，成像时的太阳高度角与光照条件相对一致，不同时相影像之间的光照差异被显著压缩，有利于开展多时相对比与变化检测——这一点对双时相模型尤为重要，因为光照与物候差异是造成伪变化的常见来源。

在载荷配置上，卫星搭载蓝、绿、红、近红外四波段多光谱传感器。多光谱信息是地物属性反演的基础：可见光波段反映地表反射特性与植被绿度状况，近红外波段对植被长势与含水量高度敏感。依托这些波段，可通过光谱特征反演土壤有机质含量、作物长势、土壤湿度等关键指标，并结合植被指数等成熟方法增强对地表覆盖状态的判别能力，从而为黑土地质量评估提供多维度数据，使变化检测结果不仅有“变了多少”的空间信息，也具备“为什么变、变得好不好”的属性解释空间。

在具体的地力指标反演层面，四个波段可组合出多种成熟的光谱指数与反演思路：以近红外与红光构建的植被指数可表征作物长势与植被覆盖度，进而反映耕地利用强度与地表裸露程度的变化；可见光波段组合有助于区分裸土、秸秆残茬与绿色植被，为保护性耕作中秸秆覆盖状况的判读提供依据；近红外波段对水分吸收的敏感性，可用于推断土壤表层湿度与作物受旱程度。将光谱反演结果与变化检测得到的空间图斑叠加分析，可把“哪里发生了变化”与“变成了什么状态”对应起来，为有机质下降、水土流失等退化判断提供更完整的证据链。

\subsubsection{（4）成熟度与可靠性：从“用得上”到“用得稳”}

该卫星技术路线获国家重点研发计划支持，已在国土资源调查、农业遥感、生态监测等领域得到广泛应用与官方认可，数据产品成熟度高、稳定性强。对项目而言，选择成熟卫星数据源的意义在于降低不确定性：数据格式规范、几何与辐射处理流程完善、产品供给稳定，可以使研发资源更多集中于变化检测算法与业务系统本身，而无需在数据获取与预处理环节反复适配。稳定的数据供给也保证了长时序监测档案的连续性，为本项目2026年落地试点、2028年推广至黑吉辽蒙四省区乃至2030年构建“天空地”一体化监测网的分阶段目标，提供了可靠的技术保障。

\subsubsection{（5）数据自主可控：从“用得上”到“用得久”}

选择国产卫星数据源，还基于数据安全与供应链自主可控的现实考量。高分辨率遥感影像的采集、存储、传输与处理涉及大量地理空间信息，其获取渠道、交付周期与使用条件若受制于外部因素，将影响监测工作的连续性与长时序数据口径的一致。采用自主可控的国产卫星数据，一方面可避免采购与授权环节的不确定性，使数据获取节奏与项目推进节奏相匹配；另一方面，国产数据产品的格式规范、坐标基准与处理工具链在国内具有更好的适配基础，便于与既有地理信息平台对接，也便于在本地化部署环境中完成数据处理，降低对外部服务依赖带来的运行风险。对面向黑土地保护的长期监测任务而言，这种稳定、可持续的数据保障能力，与算法能力和系统能力同等重要。
